%! The Rational Parent Function \& Transformations — Unit 4, Lesson 4.4 — Homework (Answer Key)
\documentclass[10pt]{article}
\usepackage{algebra2-article}
\usepackage{algebra2-key}   % pulls in -boxes; do NOT also load -boxes
\newcommand{\TallMath}[1]{$\displaystyle #1\rule[-1.4em]{0pt}{3.2em}$}
% Standard coordinate grid + axes: {xmin}{xmax}{ymin}{ymax}
\newcommand{\lgrid}[4]{%
  \draw[step=1, linegray!40, very thin] (#1,#3) grid (#2,#4);
  \draw[->, semithick, charcoal] (#1,0)--(#2,0) node[below right, font=\scriptsize]{$x$};
  \draw[->, semithick, charcoal] (0,#3)--(0,#4) node[left, font=\scriptsize]{$y$};}

\begin{document}

\pageheader{Unit 4, Lesson 4.4}{Homework --- Answer Key}
\namedateperiod

\vspace{0.10in}

\begin{notesbox}{Practice}
Every graphing question starts the same way: \textbf{find the two asymptotes.} $x=h$ comes from the
denominator (and does the opposite of what it says), $y=k$ comes from the number added outside.

\begin{enumerate}[leftmargin=1.9em, itemsep=11pt]
  \item \textbf{Asymptotes at a glance.} Give both asymptotes and describe the move from $y=\frac1x$.
    \begin{enumerate}[label=(\alph*), leftmargin=2.1em, itemsep=5pt]
      \item $y=\dfrac{1}{x-7}$ \hfill VA \ans{$x=7$} \; HA \ans{$y=0$} \; move: \ans{right $7$}
      \item $y=\dfrac{1}{x}+9$ \hfill VA \ans{$x=0$} \; HA \ans{$y=9$} \; move: \ans{up $9$}
      \item $y=-\dfrac{1}{x}$ \hfill VA \ans{$x=0$} \; HA \ans{$y=0$} \; move: \ans{reflected}
      \item $y=\dfrac{5}{x+1}$ \hfill VA \ans{$x=-1$} \; HA \ans{$y=0$} \; move: \ans{left $1$, $\times5$}
      \item $y=\dfrac{1}{x-2}+6$ \hfill VA \ans{$x=2$} \; HA \ans{$y=6$} \; move: \ans{right $2$, up $6$}
      \item $y=\dfrac{-4}{x+3}-2$ \hfill VA \ans{$x=-3$} \; HA \ans{$y=-2$} \; move: \ans{left $3$, down $2$, $\times4$, flipped}
    \end{enumerate}

  \item \textbf{The full read.} For each function give the center, both asymptotes, the domain, the
        range, and both intercepts. One of the three has no $x$-intercept --- say which and why.
    \begin{enumerate}[label=(\alph*), leftmargin=2.1em, itemsep=8pt]
      \item $g(x)=\dfrac{4}{x-1}-2$ \\[3pt]
            Center $($\ans{$1$}$,\,$\ans{$-2$}$)$ \; VA \ans{$x=1$} \; HA \ans{$y=-2$} \;
            Domain \ans{$x\neq1$} \; Range \ans{$y\neq-2$} \\[3pt]
            $y$-int \ans{$(0,-6)$} \quad $x$-int \ans{$(3,0)$}
      \item $p(x)=\dfrac{-6}{x+3}+3$ \\[3pt]
            Center $($\ans{$-3$}$,\,$\ans{$3$}$)$ \; VA \ans{$x=-3$} \; HA \ans{$y=3$} \;
            Domain \ans{$x\neq-3$} \; Range \ans{$y\neq3$} \\[3pt]
            $y$-int \ans{$(0,1)$} \quad $x$-int \ans{$(-1,0)$}
      \item $q(x)=\dfrac{1}{x+5}$ \\[3pt]
            Center $($\ans{$-5$}$,\,$\ans{$0$}$)$ \; VA \ans{$x=-5$} \; HA \ans{$y=0$} \;
            Domain \ans{$x\neq-5$} \; Range \ans{$y\neq0$} \\[3pt]
            $y$-int \ans{$(0,\frac15)$} \quad $x$-int \ans{none}
    \end{enumerate}
    \vspace{2pt}
    Which one has no $x$-intercept, and what is it about its equation that guarantees that?
    \ans{(c) --- its $k$ is $0$, so its horizontal asymptote \emph{is} the $x$-axis, and $\frac{1}{x+5}=0$ is impossible}

  \item \textbf{Write the equation from the graph.} Asymptotes are dashed; two points are marked on
        each.
    \begin{center}
    \begin{tikzpicture}[scale=0.33]
      \begin{scope}
        \lgrid{-2}{8}{-2}{6}
        \draw[navy, dashed, semithick] (3,-2)--(3,6);
        \draw[navy, dashed, semithick] (-2,2)--(8,2);
        \node[navy, font=\scriptsize, above right] at (3,5.2) {$x=3$};
        \node[navy, font=\scriptsize, below right] at (6.4,2) {$y=2$};
        \begin{scope} \clip (-2,-2) rectangle (8,6);
          \draw[very thick, forest, domain=-2:2.75, samples=110, smooth] plot (\x,{-1/((\x)-3)+2});
          \draw[very thick, forest, domain=3.25:8, samples=110, smooth] plot (\x,{-1/((\x)-3)+2});
        \end{scope}
        \fill[charcoal] (2,3) circle (3.2pt) node[above left, font=\scriptsize]{$(2,3)$};
        \fill[charcoal] (4,1) circle (3.2pt) node[below right, font=\scriptsize]{$(4,1)$};
        \node[charcoal, font=\small\bfseries] at (3,-4.2){(a)};
      \end{scope}
      \begin{scope}[xshift=15cm]
        \lgrid{-7}{5}{-5}{5}
        \draw[navy, dashed, semithick] (-2,-5)--(-2,5);
        \node[navy, font=\scriptsize, above right] at (-2,4.2) {$x=-2$};
        \begin{scope} \clip (-7,-5) rectangle (5,5);
          \draw[very thick, forest, domain=-7:-2.4, samples=110, smooth] plot (\x,{-2/((\x)+2)});
          \draw[very thick, forest, domain=-1.6:5, samples=110, smooth] plot (\x,{-2/((\x)+2)});
        \end{scope}
        \fill[charcoal] (0,-1) circle (3.2pt) node[below right, font=\scriptsize]{$(0,-1)$};
        \fill[charcoal] (-3,2) circle (3.2pt) node[above left, font=\scriptsize]{$(-3,2)$};
        \node[charcoal, font=\small\bfseries] at (-1,-7.2){(b)};
      \end{scope}
    \end{tikzpicture}
    \end{center}
    \vspace{-4pt}
    (a) $h=$ \ans{$3$} \; $k=$ \ans{$2$} \; $a=$ \ans{$-1$} \; Equation: \ans{$y=\frac{-1}{x-3}+2$}
    \\[4pt]
    (b) The horizontal asymptote is not drawn --- read it off the graph: \ans{$y=0$}. \\[3pt]
    \phantom{(b)} $h=$ \ans{$-2$} \; $k=$ \ans{$0$} \; $a=$ \ans{$-2$} \; Equation:
    \ans{$y=\frac{-2}{x+2}$}

  \item \textbf{Unmask the disguise.} Each of these \emph{is} a transformed $\frac1x$. Rewrite the
        numerator using the denominator, then read off both asymptotes.
    \begin{enumerate}[label=(\alph*), leftmargin=2.1em, itemsep=7pt]
      \item $y=\dfrac{x+5}{x+2}=\;$\ans{$1+\frac{3}{x+2}$} \quad VA \ans{$x=-2$} \quad HA \ans{$y=1$}
      \item $y=\dfrac{4x-3}{x-1}=\;$\ans{$4+\frac{1}{x-1}$} \quad VA \ans{$x=1$} \quad HA \ans{$y=4$}
      \item $y=\dfrac{2x}{x+4}=\;$\ans{$2-\frac{8}{x+4}$} \quad VA \ans{$x=-4$} \quad HA \ans{$y=2$}
    \end{enumerate}
    \vspace{2pt}
    Check all three horizontal asymptotes a second way, using Lesson 4.0's degree comparison. Do they
    agree? \ans{yes} \; Why must they? \ans{after the rewrite the leftover constant \emph{is} the ratio of the leading coefficients}

  \item \textbf{Compare and contrast.} Two tables of values:
    \begin{center}
    \renewcommand{\arraystretch}{1.3}
    \begin{minipage}{0.44\linewidth}\centering
    \textbf{Table I} \\[2pt]
    \begin{tabular}{|c|c|c|c|c|}
    \hline
    $x$ & $1$ & $2$ & $4$ & $10$ \\ \hline
    $y$ & $12$ & $6$ & $3$ & $1.2$ \\ \hline
    \end{tabular}
    \end{minipage}%
    \begin{minipage}{0.44\linewidth}\centering
    \textbf{Table II} \\[2pt]
    \begin{tabular}{|c|c|c|c|c|}
    \hline
    $x$ & $1$ & $2$ & $3$ & $4$ \\ \hline
    $y$ & $3$ & $6$ & $9$ & $12$ \\ \hline
    \end{tabular}
    \end{minipage}
    \end{center}
    \begin{enumerate}[label=(\alph*), leftmargin=2.1em, itemsep=5pt]
      \item In Table II, the \emph{differences} between consecutive $y$-values are constant, so it
            comes from a \ans{linear} function: $y=$ \ans{$3x$}.
      \item In Table I, the differences are not constant --- but every \emph{product} $x\cdot y$ is
            \ans{$12$}. So $y=$ \ans{$\frac{12}{x}$}, which is the rational parent stretched by
            \ans{$12$}.
      \item Give two features the graph of Table I's function has that Table II's does not.
            \ansline{asymptotes ($x=0$, $y=0$), a restricted domain, two branches, no intercepts --- the line has none of these}
    \end{enumerate}

  \item \textbf{Free throws.} A player has made $12$ of her first $20$ attempts. She then makes the
        next $x$ shots in a row, so her season percentage is
        \[ P(x)=\frac{12+x}{20+x}. \]
    \begin{enumerate}[label=(\alph*), leftmargin=2.1em, itemsep=5pt]
      \item $P(0)=$ \ans{$0.6$} \quad $P(20)=$ \ans{$0.8$} \quad $P(80)=$ \ans{$0.92$}
      \item Rewrite $P$ in the form $\dfrac{a}{x-h}+k$ by writing $12+x$ as $(x+20)-$\ans{$8$}:
            \\[3pt]
            $P(x)=$ \ans{$1-\frac{8}{x+20}$}, \quad so $a=$ \ans{$-8$}, \; $h=$ \ans{$-20$}, \;
            $k=$ \ans{$1$}.
      \item The horizontal asymptote is \ans{$y=1$}. What does it mean about this player?
            \ansline{her percentage climbs toward $100\%$ but never reaches it --- the $8$ misses stay in the record forever}
      \item How many in a row does she need to reach $90\%$? \ans{$60$} \; To reach $95\%$?
            \ans{$140$} \; To reach $100\%$? \ans{impossible --- no value of $x$ works}
      \item The algebra puts a vertical asymptote at $x=$ \ans{$-20$}. Which values of $x$ actually
            make sense here? \ans{whole numbers $x\ge0$}
    \end{enumerate}
\end{enumerate}
\end{notesbox}

\begin{extensionbox}
\textbf{Part 1 --- Three students, one function.} Each described $y=\dfrac{-2}{x-5}+4$. Mark each
description right or wrong, and name the error.
\begin{center}
\renewcommand{\arraystretch}{1.4}
\begin{tabularx}{\linewidth}{@{}l l X@{}}
\hline
\textbf{Student} & \textbf{Description} & \textbf{Right or wrong --- and why?} \\
\hline
Jonah & VA $x=-5$, HA $y=4$ & \ans{VA wrong --- $x-5=0$ gives $x=5$; he flipped $h$. HA correct.} \\
Kira  & VA $x=5$, HA $y=-2$ & \ans{HA wrong --- she read the numerator. $-2$ is $a$; the HA is $y=4$.} \\
Liam  & VA $x=5$, HA $y=4$; right $5$, up $4$, stretched by $2$ & \ans{Asymptotes right, description incomplete --- $a=-2$, so it is also reflected.} \\
\hline
\end{tabularx}
\end{center}
\vspace{2pt}
Liam's is the dangerous one: both asymptotes are right, and he still missed something the graph shows
plainly. What did he leave out, and where do the branches actually sit? \ansline{the reflection. With $a<0$ they sit \emph{upper-left} and \emph{lower-right} of the center $(5,4)$ --- the opposite pair from the parent.}

\vspace{4pt}
\textbf{Part 2 --- Design one.} Write the equation of a rational function whose vertical asymptote is
$x=4$, whose horizontal asymptote is $y=-3$, and whose $y$-intercept is $(0,-2)$.
\begin{center}
$y=$ \ans{$\frac{-4}{x-4}-3$}
\end{center}
\begin{enumerate}[label=(\alph*), leftmargin=1.9em, itemsep=5pt]
  \item Write your answer as a single fraction (Lesson 4.3, run forwards): \ans{$\frac{-3x+8}{x-4}$}
  \item Confirm the horizontal asymptote from that single fraction with the degree comparison.
  \item Explain why the $y$-intercept was enough to pin down $a$, but a point \emph{on} the line $y=-3$
        never could have been.
\end{enumerate}
\ansline{(b) equal degrees $\Rightarrow y=\frac{-3}{1}=-3$. (c) $(0,-2)$ gives $\frac{a}{-4}=1$; a point on $y=-3$ would force $\frac{a}{x-4}=0$, impossible --- the graph never reaches its own HA.}
\end{extensionbox}

\begin{spiralbox}
\textbf{Coming next --- Lesson 4.5: Graphing Rational Functions.} Today every function was a
transformed $\frac1x$: one factor in the denominator, one wall, asymptotes readable off the face of the
equation. Tomorrow the denominator factors into \emph{two} pieces and the picture must be assembled
instead of read --- sort each restriction into a hole or a wall (4.0), find the horizontal asymptote by
degree comparison, compute the intercepts, and use a sign chart. Everything today still applies; it
just stops being visible.
\end{spiralbox}

\begin{teachernote}
\textbf{Problems 1--2} should be quick; if not, the student is guessing $h$'s sign instead of solving
$x-h=0$. Problem 2's closing answer is \emph{$k=0$, so the HA is the $x$-axis}. \textbf{3(b)} omits the
horizontal asymptote on purpose: reading $y=0$ off the flattening branches is the task.
\textbf{Problem 4}(c) has no visible constant, so $2x=2(x+4)-8$; a student who gets $2+\frac{8}{x+4}$
should test $x=0$ ($0$ vs.\ $4$). \textbf{Problem 5} seeds \textbf{4.7} --- Table I is an inverse
variation, $xy=12$; do not name it yet. \textbf{Problem 6}(d) is worth class time: $90\%$ costs $60$
makes, $95\%$ costs $140$, $100\%$ takes forever. \textbf{Extension Part 1}: Liam is the instructive
case, since his asymptotes are right.
\end{teachernote}

\end{document}
